\section{研究ストーリー}
\label{sec:story}

\subsection{背景と問題}

製造現場ではAGV導入が進み，作業員とAGVが同一通路を共有する場面が増えている．
10〜20台規模で同時走行する環境では，
\textbf{どのAGVが自分にとって危険か}を瞬時に判断することが難しい．

経路をARで常時表示するeHMIは存在するが，
\textbf{すべてのAGVを同じ見た目で表示する}だけでは，
危険の程度を区別できず，注意資源が分散する恐れがある．

\subsection{研究課題}

\begin{quote}
  \textbf{RQ:}
  AGV残存経路の可視化において，
  \emph{危険度に応じた視覚エンコーディング}は，
  作業員の移動効率と安全性にどのような影響を与えるか？
\end{quote}

この課題に答えるため，3条件を被験者内で比較する（\Cref{tab:conditions_preview}）．

\begin{table}[htbp]
  \centering
  \caption{比較する3条件（概要）}
  \label{tab:conditions_preview}
  \begin{tabular}{lll}
    \toprule
    条件 & 経路表示 & 危険度連動 \\
    \midrule
    Baseline & あり（固定色） & なし \\
    No-AR    & なし           & --- \\
    Proposed & あり           & あり（本提案） \\
    \bottomrule
  \end{tabular}
\end{table}

\subsection{ストーリーライン}

本研究は次の論理で構成される．

\begin{enumerate}[leftmargin=2em]
  \item \textbf{問題}: 多数AGV環境での衝突リスク予測の困難さ
  \item \textbf{限界}: 均一な経路表示では危険の優先順位付けができない
  \item \textbf{提案}: TTCと近接距離から危険度 $R$ を算出し，
        色相（連続）と不透明度（4段階）で残存経路を動的に強調
  \item \textbf{検証}: VR工場で Station A$\rightarrow$B 移動タスクを実施
  \item \textbf{期待}: Proposedは No-AR より安全に，
        Baseline より情報過多を抑えつつ適切な注意喚起を実現
\end{enumerate}

\subsection{研究仮説}

\begin{description}[leftmargin=0em,style=nextline]
  \item[\textbf{H1（安全性）}]
    Proposed $<$ No-AR（ニアミス回数）\\
    \emph{根拠}: 危険AGVの経路が視覚的に強調され，回避行動が促進される．

  \item[\textbf{H2（効率）}]
    Proposed $<$ Baseline（完了時間または移動距離）\\
    \emph{根拠}: 危険度連動により，非危険AGVへの注意配分が減り，移動がスムーズになる．

  \item[\textbf{H3（情報の有用性）}]
    No-AR $>$ Proposed（完了時間または移動距離）\\
    \emph{根拠}: 経路情報がないとAGVの動きを予測できず，迂回・停止が増える．
\end{description}

\subsection{研究としての妥当性（概念的根拠）}

本研究が\textbf{意味のある問い}である理由を3点示す．

\paragraph{（1）理論的一貫性}
危険度の時間予測に TTC（Time-to-Conflict）を用いる設計は，
交通・HRI分野の慣行に沿う．
加えて近接距離を統合することで，
「交差しないが近い」「交差するがまだ遠い」両方をカバーできる．

\paragraph{（2）操作可能な独立変数}
Baseline（情報あり・均一），No-AR（情報なし），Proposed（情報あり・危険度連動）は，
\textbf{危険度エンコーディングの有無}を明確に分離する3水準である．
これにより，Proposedの効果を単独要因として解釈しやすい．

\paragraph{（3）測定可能な従属変数}
移動効率（完了時間，移動距離）と安全性（ニアミス回数）は，
VR上で自動計測可能であり，主観評価のみに依存しない．
H1--H3はいずれも，これらの客観指標で検証できる．
